% Part: many-valued-logic
% Chapter: sequent-calculus
% Section: introduction

\documentclass[../../../include/open-logic-section]{subfiles}

\begin{document}

\olfileid[pt]{mvl}{seq}{int}

\olsection{Introdução}

O cálculo de sequentes para a lógica clássica é um sistema de
!!{derivation} simples e eficiente. Se uma lógica multivalorada é
definida por uma matriz com um número finito de valores de verdade,
isto é, se $V$ é finito, é possível fornecer um cálculo de sequentes
para ela. A ideia de como fazê-lo vem da consideração dos significados
dos sequentes e da forma das regras de inferência no caso clássico.

Recorde agora que um sequente
\begin{align*}
    !A_1, \dots, !A_n & \Sequent !B_1, \dots, !B_n
\intertext{pode ser interpretado como a !!{formula}}
    (!A_1 \land \cdots \land !A_m) & \lif (!B_1 \lor \cdots \lor
!B_n)
\end{align*}
Em outras palavras, !!^a{valuation}~$\pAssign{v}$ \emph{satisfaz} um
sequente $\Gamma \Sequent \Delta$ se e somente se ou
$\pValue{v}(!A) = \False$ para algum $!A \in \Gamma$, ou
$\pValue{v}(!A) = \True$ para algum $!A \in \Delta$. Nessa
interpretação, os sequentes iniciais $!A \Sequent !A$ são sempre
satisfeitos, pois ou $\pValue v(!A) = \True$, ou
$\pValue(!A) = \False$.

Eis as regras de inferência para o condicional em~$\Log{LK}$, com as
fórmulas laterais $\Gamma$, $\Delta$ omitidas:

\begin{defish}
    \Axiom$ \fCenter !A$
    \Axiom$ !B \fCenter $
    \RightLabel{\LeftR{\lif}}
    \BinaryInf$ !A \lif !B \fCenter $
    \DisplayProof
    \hfill
    \Axiom$ !A \fCenter !B$
    \RightLabel{\RightR{\lif}}
    \UnaryInf$ \fCenter !A \lif !B $
    \DisplayProof
\end{defish}

Se aplicarmos a interpretação semântica acima a um sequente, podemos
ler a regra $\LeftR{\lif}$ como dizendo que, se
$\pValue v(!A) = \True$ e $\pValue v(!B) = \False$, então
$\pValue v(!A \lif !B) = \False$. De modo semelhante, a regra
$\RightR{\lif}$ diz que, se ou $\pValue v(!A) = \False$, ou
$\pValue v(!B) = \True$, então $\pValue v(!A \lif !B) = \True$. E, de
fato, esses condicionais são bicondicionais. No caso das regras
$\LeftR{\land}$ e $\RightR{\lor}$ em sua formulação padrão, os
condicionais correspondentes não seriam bicondicionais. Mas há versões
alternativas dessas regras nas quais eles o são:

\begin{defish}
    \Axiom$!A, !B, \Gamma \fCenter \Delta$
    \RightLabel{\LeftR{\land}}
    \UnaryInf$!A \land !B, \Gamma \fCenter \Delta$
    \DisplayProof
    \hfill
    \Axiom$ \Gamma \fCenter \Delta, !A, !B$
    \RightLabel{\RightR{\lor}}
    \UnaryInf$ \Gamma \fCenter \Delta, !A \lor !B$
    \DisplayProof
\end{defish}

Essa ideia básica, aplicada a uma lógica $n$-valorada, resulta em um
cálculo de sequentes com $n$ lugares em vez de dois, um para cada valor
de verdade. Para uma lógica trivalorada com
$V = \{\False, \Undef, \True\}$, um sequente é uma expressão
$\Gamma \mid \Pi \mid \Delta$. Ele é satisfeito em !!a{valuation}
$\pAssign v$ se e somente se ou $\pValue{v}(!A) = \False$ para algum
$!A \in \Gamma$, ou $\pValue{v}(!A) = \True$ para algum
$!A \in \Delta$, ou $\pValue{v}(!A) = \Undef$ para algum
$!A \in \Pi$. Consequentemente, os sequentes iniciais
$!A \mid !A \mid !A$ são sempre satisfeitos.

\end{document}
